\section[Human Evaluation Data Collection]{Human Evaluation Data Collection \hyperref[chsevenrq2]{\chsevenrqtwo}}\label{sec:ch-7-human-eval-setup}

We conduct a human evaluation to assess whether conditioning summary
generation on a reader's self-reported background yields summaries that
are preferred by that reader. Crucially, we measure preference with
respect to the specific annotator who provided the query and persona,
rather than aggregating judgments across annotators.

\paragraph{Annotation protocol.}
Each annotator interacts with a custom web-based platform
(see~\Cref{fig:annotation-interface} for screenshots). A session proceeds in four stages:
(i)~the annotator completes a short pre-task questionnaire eliciting
self-reported expertise, research area, and use-case, which together
constitute their \emph{persona profile};
(ii)~they submit a free-form natural-language query reflecting a
genuine information need from their own work;
(iii)~the system retrieves the top-10 paper abstracts matching the
query from OpenAlex\footnote{\url{https://openalex.org/}} and conditions generation on this
context; and
(iv)~the annotator ranks the resulting four summaries through a
pairwise tournament (described below). Each annotator completes this
loop for up to 8 distinct queries.

\paragraph{Summary generation.} For each query, four candidate summaries are produced from a
$2 \times 2$ design crossing two open-weight instruction-tuned models
Llama-3.3-70B~\cite{grattafiori2024llama3}
and DeepSeek-V3.1\footnote{\texttt{deepseek-ai/DeepSeek-V3.1}}~\cite{deepseekai2024deepseekv3technicalreport} with
two prompting conditions (with and without the persona profile). All
four summaries are generated from the same retrieved context using a
shared prompt template that instructs the model to ground its answer
in the provided abstracts and to use bracketed in-text
citations~\cite{lewis2020rag,gao2023ragsurvey}. Full prompts are
provided in~\Cref{sec:ch-7-prompts}.

\paragraph{Ranking protocol.}
Annotators rank the four summaries through a pairwise
single-elimination tournament: two pairs are compared in the first
round, and the winners advance to a final. Assignment of summaries to
bracket positions is randomized, so annotators cannot infer which
model or prompt condition produced a given summary. Pairwise
comparison has been shown to yield more reliable human judgments than
absolute Likert ratings~\cite{kiritchenko2017bws,liusie2024llmcomparative},
and the bracket structure reduces annotator effort from
$\binom{4}{2}=6$ pairwise judgments to three while producing a general
ordering over the top two summaries. For each round, the annotators also have the option to select ``Neither'' if neither summary meets their information needs. If ``Neither'' is selected in the first round bracket, the final round is skipped.

\paragraph{Annotator recruitment.}
Standard inter-annotator agreement is not a meaningful quality signal
in our setting; information satisfaction is defined relative to a
specific reader's background and query, and two annotators with
different backgrounds should be expected to disagree on the same
summary. Agreement among annotators would in fact suggest that our
protocol fails to capture the persona-conditioned signal we aim to
measure. We instead ensure annotation quality through careful
recruitment. Rather than crowdsourcing, which would require
annotators to fabricate information needs, we recruit 20 volunteer
annotators directly through professional networks, all of whom hold
genuine interest in supporting language-technology research. Our annotator pool spans domains including computer science, medicine, and materials
science, and roles including researchers, practicing professionals
(e.g., nurses), and undergraduate students. Each annotator is not required to use all assigned queries if they did not have a genuine information need for all 8. Because summaries are
generated live in response to each annotator's own query, every
judgment reflects an authentic information need. See~\Cref{sec:ch-7-annotator-comp} for details about annotator compensation.


\begin{figure*}[htbp]
    \centering
    \begin{subfigure}[t]{0.3\textwidth}
        \centering
        \includegraphics[width=\linewidth]{research-chapters/7-chapter/figures/interface-intro.png}
        \caption{Introduction to the task and instructions.}
        \label{fig:interface-intro}
    \end{subfigure}
    \hfill
    \begin{subfigure}[t]{0.2\textwidth}
        \centering
        \includegraphics[width=\linewidth]{research-chapters/7-chapter/figures/interface-profile.png}
        \caption{Profile questionnaire for building a persona.}
        \label{fig:interface-profile}
    \end{subfigure}
    \hfill
    \begin{subfigure}[t]{0.3\textwidth}
        \centering
        \includegraphics[width=\linewidth]{research-chapters/7-chapter/figures/interface-query.png}
        \caption{Annotators choose query topic.}
        \label{fig:interface-query}
    \end{subfigure}
        \vspace{1em}

    \begin{subfigure}[t]{\textwidth}
        \centering
        \includegraphics[width=.7\linewidth]{research-chapters/7-chapter/figures/interface-comparison.png}
        \caption{Example summary comparison.}
        \label{fig:interface-comparison}
    \end{subfigure}

    \caption[Annotation workflow example.]{Annotation workflow example. Steps~\ref{fig:interface-query}-\ref{fig:interface-comparison} continue for each query assigned to the annotator.}
    \label{fig:annotation-interface}
\end{figure*}